\scanpage{95}
\mspage{111}
In the case where $W=V$, $f_1:V\longrightarrow V$ is the identity and $f_o$ is the constant map, (7) reduces to the formula required in Appendix I.

The above considerations may be applied to obtain a homotopy formula for $d_\tau$, $\tau\in Y$ (see §5.2) analogous to (7) provided that $U,V,W$ and the map $f$ satisfy the additional requirement imposed by the structure and provided in addition that $U$ is pseudo-Kählerian (for example, if $U,V,W$ and the map $f$ are complex-analytic, $U$ Kählerian). The construction is based on the operator $M_{\tau,o}$, $\tau\in Y$. If the operator $M_{\tau,o}$ exists for $\tau=\sqrt{-1}$, the above argument also applies but with the difference that the term $(F_n)_U\varphi/c$ no longer vanishes. In this case $(F_n)_U=F_U$ signifies projection onto functions which are holomorphic in their dependence on the point of $U$ and we obtain the formula
\[
(\bar\partial\lambda+\lambda\bar\partial)\varphi
=f_1^*\varphi-f_o^*\varphi-((F_Uf^*)_1-(F_uf^*)_o)\varphi.
\]

The necessity of the additional terms involving $F_U$ is shown by examples (see [14]).
\mspage{94}
